\chapter{Reference Guide}
\label{toc:reference}

This chapter provides a detailed description of the side conditions generated by
the \Region plug-in, via its command-line or via its programmatic interface,
together with the associated optimizations and relevant command-line options. It
is organized by source of side-conditions: Section~\ref{toc:code} for
expression, l-values and \C instructions and Section~\ref{toc:functions} for
function pre-conditions. Special cases modify the generated side-conditions as
follows:
\begin{description}
  \item[skipped:] the annotation is not generated since it is statically verified;
  \item[invalid:] the annotation is assigned a False-if-reachable status;
\end{description}

Annotations generated through the command-line are tagged with \verb"region"
name. Invalid annotations are tagged with an additional \verb"invalid" property
name. Each category of annotations has also dedicated property names.

%% --------------------------------------------------------------------------
\section{Regions and Conditions}
\label{toc:code}
%% --------------------------------------------------------------------------

%% --------------------------------------------------------------------------
\subsection{Bounds}

\begin{lstlisting}
//@ check region: 0 <= k < N;
return a[k];
\end{lstlisting}

\noindent
Generated when accessing an array element, unless option \verb+-unsafe-arrays+
is set.

%% --------------------------------------------------------------------------
\subsection{Valid Region}

\begin{lstlisting}
//@ check region: path: \valid_region(p+k, sizeof(int));
return *(p+k);
\end{lstlisting}

\noindent
Generated when accessing a region with pointer arithmetic. Also generated when
accessing array elements when option \verb+-unsafe-arrays+ is set.

%% --------------------------------------------------------------------------
\subsection{Separated Objects}

\begin{lstlisting}
//@ check region: f: A: B: \separated(a,b);
r = f(a,b);
\end{lstlisting}

\noindent
Generated when two distinct objects \verb+A+ and \verb+B+ accessible from function
\verb+f+ might be aliased at a call site.

%% --------------------------------------------------------------------------
\subsection{Objects Attributes}

\begin{lstlisting}
//@ check region: f: A: <condition> ;
r = f(a);
\end{lstlisting}

\noindent
Generated at call site when function \verb+f+ declares an object \verb+A+. The
condition depends on the declared object attributes. However, the predicate
generated for the condition actually \emph{depends} on the memory map of each
call site: the attributes of the object \verb+A+ (from the callee specification)
are compared with the attributes of its region at call site (from the caller
memory map) to determine the residual condition to check, if any.

The general shape of condition is defined with respect of object's attributes as
specified by the table in Figure~\ref{fig:object-attributes} as combination of
object validity, nullability and initialization. Those fine-grained conditions
are generated as specified in Figure~\ref{fig:callsite-attributes} with respect
to attributes of the call site region call, here given for an object
specification \verb+A: a[p..q]+ at call site.

\begin{figure}[htbp]
\begin{center}
  \begin{tabular}{|c|l|}
    \hline
    \textbf{Object}   & \textbf{Condition} (\verb+<condition>+)\\
    \hline
    \verb+\allocated+ & \verb+<valid-read> ==> <initialized>+ \\
    \verb+\nullable+  & \verb+<non-null>   ==> <accessible>+ \\
    default           & \verb+                 <accessible>+ \\
    \hline
    \textbf{Object}   & \textbf{Accessibility} (\verb+<accessible>+)\\
    \hline
    \verb+\garbage    & \validread+   & \verb+<valid-read>+ \\
    \verb+\garbage+   & \verb+<valid-write>+ \\
    \verb+\validread+ & \verb+<valid-read>  && <initialized>+ \\
    default           & \verb+<valid-write> && <initialized>+ \\
    \hline
  \end{tabular}
\end{center}
\caption{Side Conditions for Object Attributes}
\label{fig:object-attributes}
\end{figure}

\begin{figure}[htbp]
\begin{center}
   \begin{tabular}{|c|l|}
     \hline
     \textbf{Call Site} & \textbf{Nullability} (\verb+<non-null>+)\\
     \hline
     \verb+\allocated+  & \verb$a != \null$ \\
     \verb+\nullable+   & \verb$a != \null$ \\
     default            & none (\verb+\true+) \\
     \hline
     \textbf{Call Site} & \textbf{Read Access} (\verb+<valid-read>+)\\
     \hline
     \verb+\allocated+  & \verb$\valid_read(a + (p..q))$ \\
     \verb+\nullable+   & \verb$a != \null$ \\
     default            & none (\verb+\true+) \\
     \hline
     \textbf{Call Site} & \textbf{Write Access} (\verb+<valid-write>+)\\
     \hline
     \verb+\validread+  & \verb$\valid(a + (p..q))$ \\
     \verb+\allocated+  & \verb$\valid(a + (p..q))$ \\
     \verb+\nullable+   & \verb$a != \null$ \\
     default            & none (\verb+\true+) \\
     \hline
     \textbf{Call Site} & \textbf{Initialization} (\verb+<initialized>+)\\
     \hline
     \verb+\garbage+    & \verb$\initialized(a + (p..q))$ \\
     default            & none (\verb+\true+) \\
     \hline
   \end{tabular}
\end{center}
\caption{Side Conditions at Call Site}
\label{fig:callsite-attributes}
\end{figure}

%% --------------------------------------------------------------------------
